\begin{frame}{Reformas Institucionales}
\begin{itemize}
    \item La activación de la cláusula de escape en 2025 revela los límites del marco fiscal vigente.
    \item El problema estructural es la obligación de incluir todas las obligaciones legales de gasto sin contar con ingresos suficientes.
    \item La sobreestimación de ingresos se ha convertido en una práctica recurrente para cuadrar el presupuesto.
\end{itemize}
\end{frame}

\begin{frame}{Reformas Institucionales}
\begin{itemize}
    \item Fortalecer la Oficina de Asistencia Técnica Presupuestal para aumentar la capacidad técnica del Congreso y complementar al CARF.
    \item Dar carácter independiente y vinculante a la validación de las proyecciones de ingresos y de las obligaciones legales de gasto.
    \item Si la regla fiscal impide presupuestar todas las obligaciones legales de gasto, activar un mecanismo que obligue a definir ingresos o recortes compensatorios.
    \item Modificar el artículo 348 de la Constitución para que, si el Congreso no aprueba el PGN, se reconduzca el presupuesto anterior.
\end{itemize}
\end{frame}

\begin{frame}{Reformas Institucionales}
\begin{itemize}
    \item Unificar los presupuestos de funcionamiento e inversión en el Ministerio de Hacienda.
    \item Medir la ejecución presupuestal sin contar como ejecución las transferencias a fiducias y patrimonios autónomos.
    \item Reformar las pérdidas de apropiación para permitir el traslado parcial de saldos no comprometidos.
    \item Dejar la clasificación de transacciones de única vez a cargo exclusivo del CARF.
    \item Garantizar acceso público y auditable al detalle del cierre fiscal.
\end{itemize}
\end{frame}